\documentclass[11pt]{article}
\usepackage[margin=1in]{geometry}
\usepackage{amsmath,amssymb}
\usepackage{booktabs}
\usepackage{hyperref}
\usepackage{enumitem}
\title{CCC / ADA Twin-Engine Framework}
\author{José Mancilla\\DREAMengin Research, Independent}
\date{May 2026}
\begin{document}
\maketitle

\begin{abstract}
We introduce the CCC / ADA Twin-Engine framework, a two-engine account of conserved causal compression and adaptive decoherence accounting. CCC describes what remains invariant when a physical or informational system compresses state across scale. ADA describes the participation costs, recovery ceilings, and boundary losses introduced when that compressed state is measured, decoded, or reconstructed. The central claim is that predictive structure appears only when both engines are evaluated together: conservation without accounting overstates recoverability, while accounting without conservation loses the invariant that makes prediction possible.
\end{abstract}

\section{Motivation}
Many physical descriptions separate conservation laws from measurement limits. That split is useful mathematically, but it hides a practical question: when a system compresses history into an observable state, which part of the history is conserved and which part is only recoverable inside a participation budget? The Twin-Engine answer is to treat conservation and accounting as coupled operations.

\section{Definitions}
Let $S$ be a physical state, $C(S)$ its compressed causal representation, and $A(C)$ the accounting map induced by observation or reconstruction. CCC requires an invariant $I_C$ such that
\begin{equation}
I_C(S) = I_C(C(S))
\end{equation}
for transformations inside the conserved compression class. ADA assigns a participation cost $P_A$ and a recovery ceiling $R_A$ to the same representation:
\begin{equation}
R_A = 1 - \Delta P_A.
\end{equation}
The Twin-Engine observable is therefore not $C(S)$ alone, but $(C(S), A(C(S)))$.

\section{Twin-Engine Coupling Principle}
The framework asserts that no compressed representation is physically complete until its accounting map is specified. Conversely, no accounting map is meaningful unless it acts on a conserved compressed object. This coupling is the structural bridge between gravitational residual invariance, black-hole recovery limits, and decoder performance ceilings.

\section{Predicted Signature}
The expected signature of a Twin-Engine system is scale-stable residual structure after a sign-preserving compression transform, paired with a bounded recovery or participation ceiling. This mirrors the existing Torridity record: signed-log residual structure is preserved, while the participation filter fixes a finite recovery budget.

\section{Conclusion}
CCC / ADA is not a replacement for the existing Torridity and Ledger Dynamics record. It is a broader archive frame for the same kind of result: conserved structure plus bounded participation produces falsifiable predictions.

\end{document}
